\documentclass[12pt]{article}
\usepackage[utf8]{inputenc}
\usepackage[T1]{fontenc}
\usepackage{hyperref}
\title{\textbardbl Xegwi Language Resources for K-12 Curriculum Development}
\author{Compiled by Claude Sonnet 5.5 (claude-sonnet-5-5) for Andr\'as Kornai}
\date{}
\begin{document}
\maketitle

\section{Introduction}
\textbardbl Xegwi (ISO 639-3 xeg; Glottocode xegw1238) is an extinct Tuu language of South Africa. It was spoken near Lake Chrissie, close to the Swazi border. The last known speaker is reported to have died in 1988. One news report suggests it may still be spoken in the Chrissiesmeer district, but this is unverified. Documentation is very thin. Glottolog notes that only the traditional numerals 1 and 2 were recorded. I found no educational materials, and a revival curriculum would have almost nothing to draw on.

\section{Primary Educational Resources}
None found.

\section{Linguistic and Reference Materials}
\begin{itemize}
\item \textbardbl Xegwi language (Wikipedia): \url{https://en.wikipedia.org/wiki/%C7%81Xegwi_language}
\item Tuu languages (Wikipedia): \url{https://en.wikipedia.org/wiki/Tuu_languages}
\item Glottolog entry, which cites Moseley (2010): \url{https://glottolog.org/resource/languoid/id/xegw1238}
\item Numerals page for \textbardbl Xegwi: \url{https://lingweb.eva.mpg.de/channumerals/Xegwi.htm}
\item ASJP Batwa word list, sourced from Bleek (1929), \emph{Comparative Vocabularies of Bushman Languages}. I did not verify that this list is \textbardbl Xegwi: \url{https://asjp.clld.org/languages/BATWA}
\end{itemize}

\section{Religious and Community Resources}
None found.

\section{Digital and Audio Resources}
None found.

\section{Curriculum Development Context}
Heritage work would have to be limited to what is documented: the names of the people and the language, the history of the Lake Chrissie area, and the few recorded words. Any language reconstruction would depend on specialists and on descendant communities.

\end{document}